\section{Field-Aware Few-Shot Retrieval}

DRILLER retrieves examples for schema interpretation. In a variable-schema task, semantic similarity between source texts is not enough: two passages may be topically close but ask for different fields, while two different domains may instantiate similar extraction behavior. The submitted pipelines retrieve demonstrations using schema structure, extraction instruction, field descriptions, type compatibility, and complementary output behavior. \texttt{Schema-RAG}, \texttt{Reasoned-RAG}, and every trial of \texttt{Mixed-SC} use the same retrieval path.

Figure~\ref{fig:retrieval-map} summarizes the two retrieval paths used by DRILLER.

\definecolor{drillfigblue}{RGB}{59,89,119}
\definecolor{drillfigteal}{RGB}{72,112,105}
\definecolor{drillfiggold}{RGB}{150,124,62}
\definecolor{drillfigline}{RGB}{88,88,88}

\begin{figure}[t]
\centering
\resizebox{\linewidth}{!}{%
\begin{tikzpicture}[
  font=\scriptsize,
  >=Latex,
  line cap=round,
  line join=round,
  panel/.style={
    draw=black!34,
    fill=black!2,
    rounded corners=3pt,
    line width=0.45pt
  },
  subpanel/.style={
    draw=black!28,
    fill=white,
    rounded corners=2.4pt,
    line width=0.36pt
  },
  card/.style={
    draw=black!46,
    fill=white,
    rounded corners=2pt,
    line width=0.42pt
  },
  querypart/.style={
    draw=black!25,
    fill=black!3,
    rounded corners=1.5pt,
    line width=0.32pt,
    inner sep=2.1pt,
    align=left,
    font=\tiny
  },
  badge/.style={
    draw=black!42,
    fill=white,
    rounded corners=2pt,
    line width=0.34pt,
    inner xsep=2.5pt,
    inner ysep=1.2pt,
    font=\scriptsize\bfseries
  },
  chip/.style={
    draw=black!35,
    fill=black!5,
    rounded corners=2pt,
    line width=0.30pt,
    inner xsep=2.0pt,
    inner ysep=0.9pt,
    font=\tiny
  },
  route/.style={
    -{Latex[length=2.0mm]},
    draw=drillfigline,
    line width=0.66pt
  },
  express/.style={
    -{Latex[length=2.1mm]},
    draw=drillfigblue,
    line width=1.05pt
  },
  fallback/.style={
    -{Latex[length=2.0mm]},
    draw=drillfigteal,
    line width=0.84pt,
    densely dashed
  },
  selected/.style={
    card,
    draw=drillfigblue!85!black,
    fill=drillfigblue!6,
    line width=0.82pt
  },
  selectedfallback/.style={
    card,
    draw=drillfigteal!80!black,
    fill=drillfigteal!6,
    line width=0.82pt
  },
  promptitem/.style={
    draw=black!32,
    fill=black!3,
    rounded corners=1.2pt,
    line width=0.26pt,
    inner xsep=0.85pt,
    inner ysep=0.55pt,
    outer sep=0pt,
    font=\tiny,
    align=center
  }
]

% Main grouping panels.
\node[panel, minimum width=2.62cm, minimum height=6.16cm, anchor=south west] (querypanel) at (0.05,0.62) {};
\node[panel, minimum width=6.78cm, minimum height=2.48cm, anchor=south west, fill=drillfigblue!4] (upperpanel) at (3.05,4.33) {};
\node[panel, minimum width=6.78cm, minimum height=3.28cm, anchor=south west, fill=drillfigteal!4] (lowerpanel) at (3.05,0.62) {};
\node[panel, minimum width=2.58cm, minimum height=6.16cm, anchor=south west] (selectedpanel) at (10.10,0.62) {};
\node[panel, minimum width=3.07cm, minimum height=6.16cm, anchor=south west] (promptpanel) at (12.68,0.62) {};

\node[badge, fill=black!4, anchor=west, inner xsep=1.6pt] at (0.24,6.55) {Query instance \(q\)};
\node[badge, draw=drillfigblue!60!black, fill=drillfigblue!8, anchor=west] at (3.24,6.56) {same-schema retrieval};
\node[chip, draw=drillfigblue!65!black, fill=white, anchor=east] at (9.62,6.55) {FAST PATH};
\node[badge, draw=drillfigteal!65!black, fill=drillfigteal!8, anchor=west] at (3.24,3.65) {field-level fallback};
\node[chip, draw=drillfigteal!65!black, fill=white, anchor=east] at (9.62,3.64) {ANALYTICAL};
\node[badge, fill=black!4, anchor=west] at (10.25,6.55) {Retrieved};
\node[badge, fill=black!4, anchor=west] at (12.83,6.55) {Prompt integration};

% Query card.
\node[card, minimum width=2.12cm, minimum height=5.22cm, anchor=north west] (querycard) at (0.36,6.23) {};
\draw[draw=black!34, fill=black!10, line width=0.32pt]
  (2.19,6.23) -- (2.48,6.23) -- (2.48,5.94) -- cycle;
\draw[draw=black!28, line width=0.32pt] (2.19,6.23) -- (2.48,5.94);
\node[querypart, text width=1.62cm, anchor=north west] at (0.53,5.88)
  {\textbf{instruction}\\extract fields};
\node[querypart, text width=1.62cm, anchor=north west] at (0.53,4.90)
  {\textbf{source text}\\short passage};
\node[querypart, text width=1.62cm, anchor=north west] at (0.53,3.92)
  {\textbf{target JSON Schema}\\types, enums, req.};
\node[font=\tiny\bfseries, text=black!65, anchor=west] at (0.53,2.69) {extracted views};
\node[chip, draw=drillfigblue!55!black, fill=drillfigblue!7, anchor=west] at (0.53,2.36) {\(\phi(S_q)\)};
\node[chip, draw=drillfigteal!60!black, fill=drillfigteal!7, anchor=west] at (0.53,2.05) {\(f_1,\ldots,f_m\)};
\node[font=\tiny, text=black!60, align=center, text width=1.72cm] at (1.38,1.58)
  {schema identity\\and field list};

% Same-schema lane.
\coordinate (qout) at (2.67,5.48);
\draw[route] (2.48,5.48) -- (qout);
\draw[express] (qout) -- (3.12,5.48);

\node[font=\tiny\bfseries, text=black!70, anchor=west] at (3.24,6.12) {FSP candidates};
\foreach \x/\lab/\shade in {3.32/\(c_1\)/4,3.69/\(c_2\)/9,4.06/\(c_3\)/4,4.35/\(\cdots\)/11} {
  \node[card, fill=black!\shade, minimum width=0.38cm, minimum height=0.68cm] at (\x,5.48) {};
  \node[font=\tiny] at (\x,5.48) {\lab};
}
\draw[express] (4.54,5.48) -- (4.98,5.48);
\node[card, draw=drillfigblue!65!black, fill=white, minimum width=0.95cm, minimum height=1.18cm, align=center, font=\scriptsize] (gate) at (5.52,5.48)
  {\(\phi(S_q)\)\\[-1pt]\(=\phi(S_c)\)};
\node[font=\tiny, text=black!62] at (5.52,4.75) {hard filter};
\draw[express] (6.01,5.48) -- (6.40,5.48);

\node[card, minimum width=1.50cm, minimum height=1.22cm] (rankbox) at (7.17,5.48) {};
\node[font=\tiny\bfseries, align=center] at (7.17,5.90) {semantic\\rank};
\node[font=\tiny] at (7.17,5.58) {\(\operatorname{sim}(r_q,r_c)\)};
\draw[fill=drillfigblue!55, draw=black!20, line width=0.22pt] (6.62,5.34) rectangle (7.48,5.43);
\draw[fill=drillfigblue!35, draw=black!20, line width=0.22pt] (6.62,5.16) rectangle (7.28,5.25);
\draw[fill=black!22, draw=black!20, line width=0.22pt] (6.62,4.98) rectangle (7.02,5.07);
\draw[express] (7.91,5.48) -- (8.24,5.48);

% Explicit same-schema decision node.
\coordinate (dec) at (8.76,5.48);
\draw[draw=drillfigblue!70!black, fill=white, line width=0.55pt]
  (8.76,5.98) -- (9.28,5.48) -- (8.76,4.98) -- (8.24,5.48) -- cycle;
\node[font=\tiny\bfseries, align=center] at (8.76,5.55) {\(\geq 2\)\\pass \(\tau\)?};
\node[font=\tiny, text=drillfigblue!70!black, anchor=south] at (9.46,5.62) {YES};
\draw[express] (9.28,5.48) -- (10.10,5.48);

% Fallback control flow from the decision.
\node[font=\tiny, text=drillfigteal!75!black, align=center, text width=0.88cm, anchor=north] at (9.38,5.16)
  {NO:\\field-level\\fallback};
\draw[fallback] (8.76,4.98) -- (8.76,4.18) -- (3.36,4.18) -- (3.36,3.90);

% Field-level fallback lane.
\node[font=\tiny\bfseries, text=black!70, anchor=west] at (3.34,3.28) {type-compatible field similarity};
\foreach \x/\lab in {3.86/\(c_1\),4.24/\(c_2\),4.62/\(c_3\),5.00/\(\cdots\)} {
  \node[font=\tiny] at (\x,3.02) {\lab};
}
\foreach \y/\lab in {2.72/\(f_1\),2.42/\(f_2\),2.12/\(f_3\),1.82/\(f_m\)} {
  \node[font=\tiny, anchor=east] at (3.62,\y) {\lab};
}
\foreach \x/\y/\shade in {
  3.70/2.58/65,4.08/2.58/28,4.46/2.58/12,4.84/2.58/45,
  3.70/2.28/18,4.08/2.28/62,4.46/2.28/34,4.84/2.28/10,
  3.70/1.98/48,4.08/1.98/14,4.46/1.98/58,4.84/1.98/22,
  3.70/1.68/10,4.08/1.68/42,4.46/1.68/26,4.84/1.68/60
} {
  \draw[draw=black!18, fill=black!\shade, line width=0.23pt] (\x,\y) rectangle ++(0.30,0.26);
}
\draw[draw=black!40, line width=0.32pt] (3.70,1.68) rectangle (5.14,2.84);
\draw[draw=drillfigteal!78!black, line width=0.58pt, rounded corners=1pt] (3.67,1.64) rectangle (4.03,2.88);
\draw[draw=drillfigteal!78!black, line width=0.58pt, rounded corners=1pt, densely dashed] (4.05,1.64) rectangle (4.41,2.88);
\node[chip, draw=drillfigteal!60!black, fill=white, font=\tiny] at (4.42,1.23)
  {\(\kappa(f_i)\sim\kappa(g)\)};
% Field-similarity vectors tied to selected heatmap columns.
\draw[fallback, line width=0.58pt] (4.03,2.64) -- (5.50,2.42);
\draw[fallback, line width=0.58pt] (4.41,2.24) -- (5.50,1.52);
\node[font=\tiny, anchor=south] at (6.02,2.82) {\(v(c_1)\)};
\node[card, minimum width=1.05cm, minimum height=0.72cm] (vonebox) at (6.02,2.38) {};
\draw[fill=drillfigteal!55, draw=black!20, line width=0.2pt] (5.70,2.14) rectangle ++(0.10,0.32);
\draw[fill=drillfigteal!25, draw=black!20, line width=0.2pt] (5.88,2.14) rectangle ++(0.10,0.15);
\draw[fill=drillfigteal!45, draw=black!20, line width=0.2pt] (6.06,2.14) rectangle ++(0.10,0.27);
\draw[fill=black!20, draw=black!20, line width=0.2pt] (6.24,2.14) rectangle ++(0.10,0.10);

\node[card, minimum width=1.05cm, minimum height=0.72cm] (vtwobox) at (6.02,1.50) {};
\draw[fill=drillfigteal!25, draw=black!20, line width=0.2pt] (5.70,1.26) rectangle ++(0.10,0.13);
\draw[fill=drillfigteal!55, draw=black!20, line width=0.2pt] (5.88,1.26) rectangle ++(0.10,0.34);
\draw[fill=black!20, draw=black!20, line width=0.2pt] (6.06,1.26) rectangle ++(0.10,0.11);
\draw[fill=drillfigteal!48, draw=black!20, line width=0.2pt] (6.24,1.26) rectangle ++(0.10,0.28);
\node[font=\tiny, anchor=north] at (6.02,1.05) {\(v(c_2)\)};

\draw[fallback] (6.58,2.02) -- (6.82,2.02);
\node[card, minimum width=1.78cm, minimum height=0.88cm, align=center, font=\tiny] (objective) at (8.05,2.80)
  {\(s(c_1,c_2)=\)\\[-1pt]
   \(\|v(c_1)+v(c_2)\|_2\)\\[-1pt]
   \(+\|v(c_1)-v(c_2)\|_2\)};
\node[card, minimum width=1.36cm, minimum height=0.80cm] (joint) at (7.35,1.48) {};
\node[font=\tiny\bfseries] at (7.35,1.78) {coverage};
\node[font=\tiny] at (7.35,1.58) {\(v_1+v_2\)};
\draw[fill=drillfiggold!55, draw=black!20, line width=0.2pt] (6.92,1.15) rectangle ++(0.14,0.22);
\draw[fill=drillfiggold!50, draw=black!20, line width=0.2pt] (7.15,1.15) rectangle ++(0.14,0.31);
\draw[fill=drillfiggold!38, draw=black!20, line width=0.2pt] (7.38,1.15) rectangle ++(0.14,0.25);
\draw[fill=drillfiggold!48, draw=black!20, line width=0.2pt] (7.61,1.15) rectangle ++(0.14,0.28);
\node[card, minimum width=1.50cm, minimum height=0.80cm] (diff) at (8.85,1.48) {};
\node[font=\tiny\bfseries] at (8.85,1.78) {complement};
\node[font=\tiny] at (8.85,1.58) {\(v_1-v_2\)};
\draw[draw=black!38, line width=0.25pt] (8.35,1.28) -- (9.36,1.28);
\draw[fill=drillfigteal!45, draw=black!20, line width=0.2pt] (8.40,1.28) rectangle ++(0.16,0.17);
\draw[fill=drillfigteal!30, draw=black!20, line width=0.2pt] (8.68,1.16) rectangle ++(0.16,0.12);
\draw[fill=drillfigteal!42, draw=black!20, line width=0.2pt] (8.96,1.28) rectangle ++(0.16,0.14);
\draw[fill=drillfigteal!35, draw=black!20, line width=0.2pt] (9.24,1.14) rectangle ++(0.16,0.14);
\draw[fallback] (9.52,2.20) -- (10.10,2.20);

% Selected demonstrations panel, split by retrieval mode.
\node[subpanel, minimum width=2.26cm, minimum height=2.51cm, anchor=south west] (sameselpanel) at (10.26,3.82) {};
\node[subpanel, minimum width=2.26cm, minimum height=2.62cm, anchor=south west] (fbselpanel) at (10.26,0.88) {};
\node[font=\tiny\bfseries, anchor=west] at (10.38,6.10) {same-schema outputs};
\node[font=\tiny, text=black!62, anchor=west] at (10.38,5.87) {shared schema};
\node[selected, text width=0.72cm, minimum height=0.70cm, align=center, font=\tiny] at (10.82,5.42)
  {\textbf{Ex. 1}\\demo};
\node[selected, text width=0.72cm, minimum height=0.70cm, align=center, font=\tiny] at (11.72,5.42)
  {\textbf{Ex. 2}\\demo};
\draw[draw=drillfigblue!65!black, line width=0.45pt] (10.50,5.05) -- (12.05,5.05);
\node[chip, draw=drillfigblue!55!black, fill=white] at (11.26,4.78) {null vs grounded};
\node[chip, draw=drillfigblue!55!black, fill=white] at (11.26,4.50) {empty vs populated};
\node[chip, draw=drillfigblue!55!black, fill=white] at (11.26,4.22) {enum contrast};

\node[font=\tiny\bfseries, anchor=west] at (10.38,3.34) {fallback outputs};
\node[font=\tiny, text=black!62, anchor=north west, align=left, text width=1.36cm, inner sep=0pt] at (10.38,3.15) {complementary\\[-1pt]coverage};
\node[selectedfallback, text width=1.74cm, minimum height=0.82cm, align=left, font=\tiny] at (11.39,2.24)
  {\textbf{Ex. 1}\hfill \(f_1\ f_3\)};
\draw[fill=drillfigteal!55, draw=black!20, line width=0.18pt] (10.62,1.90) rectangle ++(0.20,0.22);
\draw[fill=black!15, draw=black!20, line width=0.18pt] (10.91,1.90) rectangle ++(0.20,0.08);
\draw[fill=drillfigteal!48, draw=black!20, line width=0.18pt] (11.20,1.90) rectangle ++(0.20,0.19);
\draw[fill=black!15, draw=black!20, line width=0.18pt] (11.49,1.90) rectangle ++(0.20,0.07);
\node[selectedfallback, text width=1.74cm, minimum height=0.82cm, align=left, font=\tiny] at (11.39,1.36)
  {\textbf{Ex. 2}\hfill \(f_2\ f_m\)};
\draw[fill=black!15, draw=black!20, line width=0.18pt] (10.62,1.02) rectangle ++(0.20,0.08);
\draw[fill=drillfigteal!56, draw=black!20, line width=0.18pt] (10.91,1.02) rectangle ++(0.20,0.24);
\draw[fill=black!15, draw=black!20, line width=0.18pt] (11.20,1.02) rectangle ++(0.20,0.08);
\draw[fill=drillfigteal!48, draw=black!20, line width=0.18pt] (11.49,1.02) rectangle ++(0.20,0.20);
% Prompt integration panel, split by rendering layout.
\node[subpanel, minimum width=2.86cm, minimum height=2.23cm, anchor=south west] (samepromptpanel) at (12.78,4.10) {};
\node[subpanel, minimum width=2.86cm, minimum height=2.92cm, anchor=south west] (ragpromptpanel) at (12.78,0.88) {};
\node[font=\tiny\bfseries, anchor=west] at (12.90,6.13) {same-schema assembly};
\node[promptitem, minimum width=0.80cm, minimum height=0.22cm, fill=black!3] at (13.94,5.84) {rules};
\node[promptitem, minimum width=1.16cm, minimum height=0.22cm, fill=drillfigblue!7] at (14.84,5.84) {schema \(S_q\)};
\foreach \y/\lab in {5.38/Ex. 1,5.00/Ex. 2} {
  \node[font=\tiny\bfseries, anchor=east] at (13.64,\y) {\lab};
  \node[promptitem, minimum width=0.48cm, minimum height=0.22cm] at (13.96,\y) {instr};
  \node[promptitem, minimum width=0.44cm, minimum height=0.22cm] at (14.54,\y) {src};
  \node[promptitem, minimum width=0.44cm, minimum height=0.22cm] at (15.10,\y) {out};
  \draw[draw=black!38, -{Latex[length=0.8mm]}, line width=0.24pt] (14.24,\y) -- (14.30,\y);
  \draw[draw=black!38, -{Latex[length=0.8mm]}, line width=0.24pt] (14.82,\y) -- (14.88,\y);
}
\node[font=\tiny\bfseries, anchor=east] at (13.72,4.62) {current};
\node[promptitem, minimum width=0.48cm, minimum height=0.22cm, fill=drillfigblue!5] at (13.96,4.62) {instr};
\node[promptitem, minimum width=0.44cm, minimum height=0.22cm, fill=drillfigblue!5] at (14.54,4.62) {src};
\draw[draw=black!38, -{Latex[length=0.8mm]}, line width=0.24pt] (14.24,4.62) -- (14.30,4.62);

\node[font=\tiny\bfseries, anchor=west] at (12.90,3.57) {general RAG assembly};
\foreach \y/\lab/\schema in {3.08/Ex. 1/\(S_{c_1}\),2.32/Ex. 2/\(S_{c_2}\)} {
  \node[promptitem, minimum width=1.72cm, minimum height=0.56cm] at (14.56,\y) {};
  \node[font=\tiny\bfseries, anchor=east] at (13.56,\y+0.16) {\lab};
  \node[promptitem, minimum width=0.48cm, minimum height=0.18cm] at (14.06,\y+0.16) {instr};
  \node[promptitem, minimum width=0.62cm, minimum height=0.18cm, fill=drillfigteal!7] at (14.80,\y+0.16) {\schema};
  \node[promptitem, minimum width=0.46cm, minimum height=0.18cm] at (14.06,\y-0.15) {src};
  \node[promptitem, minimum width=0.46cm, minimum height=0.18cm] at (14.80,\y-0.15) {out};
  \draw[draw=black!38, -{Latex[length=0.8mm]}, line width=0.24pt] (14.34,\y-0.15) -- (14.50,\y-0.15);
}
\node[promptitem, minimum width=1.72cm, minimum height=0.48cm, fill=drillfigteal!6] at (14.56,1.52) {};
\node[font=\tiny\bfseries, anchor=east] at (13.72,1.64) {current};
\node[promptitem, minimum width=0.48cm, minimum height=0.18cm, fill=white] at (14.06,1.64) {instr};
\node[promptitem, minimum width=0.62cm, minimum height=0.18cm, fill=white] at (14.80,1.64) {\(S_q\)};
\node[promptitem, minimum width=0.46cm, minimum height=0.18cm, fill=white] at (14.44,1.32) {src};

% Mode-specific integration arrows.
\draw[express] (12.52,5.48) -- (12.82,5.48);
\draw[fallback] (12.52,2.20) -- (12.82,2.20);

\end{tikzpicture}%
}
\caption{Two-stage retrieval in DRILLER. The system first attempts same-schema retrieval by filtering candidate demonstrations with a normalized schema fingerprint and ranking surviving cases semantically. A decision node routes directly to same-schema demonstrations when at least two cases pass the threshold; otherwise, field-level fallback scores candidates by type-compatible field similarity and selects a complementary pair. The selected demonstrations are then rendered with either the compact same-schema prompt layout or the general RAG layout.}
\label{fig:retrieval-map}
\end{figure}


\subsection{Curated Demonstration Corpus}

The retrieval pool is a local few-shot prompting (FSP) corpus of structured extraction demonstrations. Each case contains a Spanish source text, an extraction instruction, a target schema, expected values, tags describing covered extraction phenomena, and field-level rationales or descriptions that support rendering examples in either direct or reasoning-augmented form. The corpus is synthetic and curated rather than mined from an external collection.

The GenSIE development set guided corpus design. We used it to study schema structure, source-text style, vocabulary, noise patterns, and recurring ambiguities, rather than reusing development instances as demonstrations. We also manually curated a small development-time evaluation subset, reviewing up to two instances per task family. Those curated instances served as design references only. Direct reuse would not provide systematic complementary coverage of field dualities and would contaminate evaluation on the same curated subset, since retrieval could return the instance being solved.

We constructed the final FSP corpus as a compact library of synthetic and manually curated structured cases, organized by development task family and schema family. The corpus used an author-in-the-loop workflow assisted by an agentic LLM coding assistant. The assistant helped inspect candidate task families, draft complementary-case proposals, generate provisional source texts and structured resources, and run structural checks. The authors reviewed and corrected the proposed dualities, source texts, expected outputs, reasoning fields, enriched descriptions, and final JSON resources before inclusion. The workflow was:

\begin{enumerate}
\item Each development task family became the unit of work because a family bundles a schema, domain, source-text style, vocabulary, and recurring extraction traps.
\item The assistant inspected the curated development subset first, then consulted the full uncurated development set when more evidence about style, structure, or ambiguity was needed. The full set characterized text length, noise patterns, and extraction behavior, but was not treated as trusted gold output.
\item The assistant identified useful field-level dualities: nullable value versus grounded non-null value; empty array versus populated array; boolean \texttt{true} versus \texttt{false}; frequent enum label versus an alternative label; direct or near-verbatim text versus grounded summary; explicit value versus distractor or insufficient evidence; and normalized value versus literal mention.
\item It then proposed two complementary cases for the task or schema family, using a balanced or approximately balanced mask over dual fields as a design heuristic. One case covered one side of several dualities and the second covered the inverse side; qualitative contrasts such as enum A versus enum B or explicit date versus nearby distractor date were handled outside a rigid binary mask.
\item The authors reviewed this pre-drafting proposal, including selected references, field dualities, complementary mask, and intended case behavior. Full synthetic drafting began only after this gate was accepted.
\item Under the approved design, the assistant drafted synthetic Spanish source texts, candidate expected outputs, and field-level reasoning. The texts had to ground values, \texttt{null}s, and empty arrays naturally, resemble the reviewed family style, and avoid benchmark-aware phrases that mechanically announced the expected answer or absence.
\item The authors manually curated the drafted cases before final resource generation, editing source texts, expected outputs, and field-level reasoning. This review checked that each value, \texttt{null}, and empty list was justified by the source text and that the pair contrasted the intended regimes.
\item After author approval, the assistant generated the final \texttt{StructuredFspCase} JSON resources, including \texttt{field\_examples}, tags, field-level reasoning, and enriched descriptions when applicable. The resources were loaded through the FSP resource loader or equivalent structural checks before retrieval used them.
\end{enumerate}

The final extraction corpus contains two cases for each of the 22 schema or task families represented in the curated development setting, for 44 extraction cases in the local retrieval pool. The families cover a range of domains and structures, but the corpus is not intended to be a large topical nearest-neighbor database. It is a compact library of contrasted demonstrations: the useful unit is often a pair of cases showing how the same or related fields behave under different evidence conditions.

The paired organization is deliberate. A nullable field benefits from both grounded non-null values and grounded absence. An array field benefits from examples where the correct list is populated and examples where it is empty. Enum fields benefit from contrasting labels, especially when labels are inferred rather than copied literally. Numeric and date fields may require exact copying in one case and normalization in another. String fields may require direct mention extraction, near-verbatim evidence, or a short grounded summary depending on the field description.

After complementary cases had been proposed and reviewed, we added enriched field descriptions as compact Spanish prompt annotations. They explain how to extract a field and warn about traps observed during development-example review or synthetic case construction, such as distractor dates, nearby but semantically wrong entities, aliases, or conditions under which a nullable field should remain \texttt{null}. These descriptions clarify schema semantics for same-schema prompting; they do not add new instance evidence. They appear only in the prompt-side schema view. The target JSON Schema, strict generation schema, and evaluator-facing output schema remain unchanged.

The FSP reasoning fields follow the structured evidence discipline. In reasoning-mode examples, each field can render \texttt{field\_asks}, \texttt{relevant\_fragments}, and \texttt{final\_value} as a local evidence protocol. Cited fragments support the extraction decision rather than merely repeating the answer string. They stay short while preserving the context needed to justify the value: accepted candidates show why the value belongs to the requested field, and rejected or \texttt{null} cases cite the closest reviewed context and state which required assertion is missing. Enum reasoning considers the allowed labels and justifies the selected label from relevant fragments. Array reasoning may cite a complete list once, or cite compactly around individual items, so examples support recall without bloating the prompt. FSP examples therefore demonstrate grounding behavior as well as output formatting.

Table~\ref{tab:technical-software-fsp-pair} gives one compact example from the \texttt{technical\_software} family. The two cases share the same schema but expose complementary regimes: one case grounds populated arrays and leaves several adversarial nullable fields unsupported, while the other inverts the pattern with empty arrays, alias resolution, exact date normalization, an explicit stable-installer checksum, and explicit corporate metadata.

\begin{table}[t]
\centering
\footnotesize
\setlength{\tabcolsep}{3pt}
\caption{Complementary FSP pair for the \texttt{technical\_software} schema family.}
\label{tab:technical-software-fsp-pair}
\begin{tabularx}{\linewidth}{@{}>{\raggedright\arraybackslash}p{0.20\linewidth}
>{\raggedright\arraybackslash}X
>{\raggedright\arraybackslash}X@{}}
\toprule
Field or contrast & \textit{Lince Editor} case & \textit{QuantaPack Pro} case \\
\midrule
\texttt{official\_name} and \texttt{app\_category} &
Direct product name; enum maps to \texttt{IDE\_EDITOR}. &
Abbreviated heading must resolve to full name; enum maps to \texttt{ARCHIVER}. \\
\texttt{platforms} and \texttt{features} &
Populated arrays with named operating systems and technical capabilities. &
Empty arrays: installers and downloads are mentioned, but no supported OS or concrete feature list. \\
\texttt{exact\_release\_date} &
\texttt{null}: only year-level or relative release evidence is present. &
Exact first-release date normalized to \texttt{DD/MM/YYYY}. \\
\path{latest_stable_version_sha256} &
\texttt{null}: PGP signatures and a nightly commit hash are distractors. &
Explicit SHA256 for the stable installer. \\
\texttt{current\_ceo\_name} &
\texttt{null}: founder, idea author, and coordinating team are not CEO evidence. &
Explicit current CEO of the developing company. \\
Pair lesson &
Extract grounded scalar values and populated arrays, but abstain for nullable fields with only distractor evidence. &
Resolve aliases, keep unsupported arrays empty, and fill nullable fields only when the source states the required evidence. \\
\bottomrule
\end{tabularx}
\end{table}

Coverage includes null versus non-null values, present versus absent evidence, enum alternatives, empty versus populated arrays, arrays of simple values, arrays of objects, nested objects, direct strings, summary strings, boolean inference, numeric normalization, date normalization, bounded scores, long evidence spans, and distractor mentions. These contrasts matter because few-shot examples can bias the model when they show only one outcome regime. Complementary demonstrations teach conditional behavior: extract when evidence exists, abstain when it does not, normalize only when requested, and respect the requested type and label set.

\subsection{Same-Schema Retrieval}

At runtime, DRILLER first tries to retrieve examples with the same normalized schema as the current task. Let \(S_q\) be the target schema of query instance \(q\), and let \(S_c\) be the schema of a candidate case \(c\). DRILLER computes a normalized schema fingerprint \(\phi(S)\) by canonicalizing the JSON representation while preserving the structural extraction contract: field names, object structure, array structure, primitive types, enum alternatives, required properties, and nesting. Descriptions are neutralized so that prompt wording refinements do not change schema identity. Order-insensitive elements such as \texttt{required}, \texttt{enum}, and array-valued \texttt{type} are sorted before serialization. Same-schema candidates satisfy the hard constraint
\[
\phi(S_q)=\phi(S_c).
\]

The fingerprint separates schema identity from topical similarity. A medical passage and a technical passage might both contain names, dates, arrays, and classifications without sharing an extraction form. Conversely, two source texts from the same schema family may require the same field interpretation even when their values differ. Same-schema retrieval therefore treats structural identity as a filter rather than as a soft ranking feature.

Filtered candidates are then ranked semantically. DRILLER builds a global textual representation \(r(q)\) of the current task from the task identifier, extraction instruction, and textual representations of top-level fields. Candidate cases are represented analogously as \(r(c)\). These representations are embedded and compared by cosine similarity,
\[
\operatorname{sim}(r(q),r(c)).
\]
The submitted system uses a same-schema threshold \(\tau=0.6\). If at least two fingerprint-matching candidates satisfy \(\operatorname{sim}(r(q),r(c))\geq\tau\), DRILLER selects the top candidates by similarity and marks them as same-schema demonstrations. If fewer than two candidates pass this test, retrieval continues with the field-level fallback.

Same-schema examples show the exact output contract under different source texts. They clarify what counts as evidence for a field, when null or an empty list is correct, and how enum labels or nested structures are populated. They also enable a compact prompt layout: the shared schema is rendered once, and examples focus on instruction, source text, and expected output rather than repeating identical schema blocks.

Same-schema mode can also enrich the prompt-side schema view with curated field descriptions drawn from the selected cases. These enriched descriptions add task and domain guidance such as enum interpretation, null traps, or distractor patterns to the schema text shown to the model. They do not modify the target JSON Schema, the strict generation schema, or the final evaluator-facing schema.

\subsection{Field-Level Retrieval Fallback}

When same-schema retrieval does not produce enough candidates, DRILLER falls back to field-level retrieval. Let the query top-level fields be \(F_q=\{f_1,\ldots,f_m\}\), in schema order, and let \(F_c\) be the top-level fields of candidate case \(c\). Each field \(f\) is mapped to a textual representation \(t(f)\) that combines its field name, readable type, and description, for example ``\texttt{name (type): description}''. The retriever also assigns a coarse type class \(\kappa(f)\), ignoring nullability.

The coarse classes used by the retriever are \texttt{string}, \texttt{boolean}, \texttt{numeric} for both integer and number fields, \texttt{enum}, \texttt{object}, and recursively typed arrays such as \texttt{array[string]}, \texttt{array[numeric]}, \texttt{array[enum]}, and \texttt{array[object]}. A fallback \texttt{any} class is available for unsupported or underspecified schemas. We write \(\kappa(f_i)\sim\kappa(g)\) for compatibility; in the submitted implementation this relation is equality of the nullability-insensitive coarse classes. Thus nullable and non-nullable variants of the same base type can match, while an array of objects is not aligned with a scalar summary merely because their descriptions share topical words.

For each candidate \(c\), DRILLER embeds the field texts \(e(t(f))\) and computes, for every query field \(f_i\), the best compatible field match in the candidate:
\[
a_i(c)=
\max\left(0,\max_{g\in F_c,\ \kappa(g)\sim\kappa(f_i)}
\cos(e(t(f_i)),e(t(g)))\right),
\]
with \(a_i(c)=0\) when no compatible candidate field exists. The candidate is therefore represented by a field-similarity vector aligned with the current query schema,
\[
v(c)=(a_1(c),\ldots,a_m(c)).
\]
A candidate's individual relevance can be summarized by \(\lVert v(c)\rVert_2\), but two-example selection uses a pair-level objective over these same vectors.

The type-aware vector prevents topical retrieval from ignoring output shape. It avoids matching semantically similar descriptions with incompatible JSON structures, allows nullable and non-nullable variants of the same base type to reinforce each other, and shows whether a candidate is broadly useful or strong for only part of the schema. Since retrieval selects a pair of demonstrations, two individually imperfect candidates can be preferable when they cover different field behaviors.

If embedding retrieval fails or produces no usable candidates, DRILLER uses a lexical/schema fallback. It scores candidates with explicit features such as exact normalized schema matches, compatible field fingerprints, tag overlap, field-name overlap, shared terms, and domain-term overlap, with a small penalty for long example texts. This fallback preserves schema-aware retrieval when embeddings are unavailable or uninformative.

\subsection{Complementary Pair Selection}

The submitted RAG path retrieves up to two examples. In the field-level fallback, DRILLER selects a pair with a complementarity-aware objective rather than choosing the two highest individual scores. For two candidate demonstrations \(c_1\) and \(c_2\), with field-similarity vectors \(v(c_1)\) and \(v(c_2)\) defined above, the pair score is

\[
s(c_1,c_2) = \lVert v(c_1)+v(c_2)\rVert_2 + \lVert v(c_1)-v(c_2)\rVert_2 .
\]

The first term, \(\lVert v(c_1)+v(c_2)\rVert_2\), rewards joint relevance and field coverage. If both examples provide useful analogues for the current schema, their combined vector has large magnitude. This term aligns with the ordinary retrieval goal that demonstrations should be relevant to the task.

The second term, \(\lVert v(c_1)-v(c_2)\rVert_2\), rewards complementarity. If two candidates are strong on the same fields and weak on the same fields, their difference vector is small. If one is strong where the other is weak, the difference grows. This discourages near-duplicate demonstrations that reinforce one behavior while leaving other fields unexplained.

\begin{figure}[t]
\centering
\resizebox{0.82\linewidth}{!}{%
\begin{tikzpicture}[
  font=\scriptsize,
  >=Latex,
  x={(1.65cm,0cm)},
  y={(-0.42cm,0.34cm)},
  z={(0cm,1.52cm)},
  axis3/.style={->, line width=0.55pt},
  cube/.style={line width=0.35pt, gray!45},
  candA/.style={->, line width=1.0pt, blue!70!black},
  candB/.style={->, line width=1.0pt, orange!85!black},
  sumvec/.style={->, dashed, line width=0.85pt, green!45!black},
  guide/.style={densely dotted, line width=0.35pt, gray!65},
  diffline/.style={<->, line width=0.5pt, gray!75!black},
  difflabel/.style={inner sep=0.6pt, gray!70!black},
  vlabel/.style={inner sep=0.6pt}
]

\begin{scope}
\node[font=\bfseries\scriptsize] at (0.48,0.50,1.48) {Redundant pair};
\fill[gray!5] (0,0,0) -- (1,0,0) -- (1,1,0) -- (0,1,0) -- cycle;
\fill[gray!9] (0,1,0) -- (1,1,0) -- (1,1,1) -- (0,1,1) -- cycle;
\fill[gray!13] (1,0,0) -- (1,1,0) -- (1,1,1) -- (1,0,1) -- cycle;
\draw[cube] (0,0,0) -- (1,0,0) -- (1,1,0) -- (0,1,0) -- cycle;
\draw[cube] (0,0,1) -- (1,0,1) -- (1,1,1) -- (0,1,1) -- cycle;
\draw[cube] (0,0,0) -- (0,0,1);
\draw[cube] (1,0,0) -- (1,0,1);
\draw[cube] (0,1,0) -- (0,1,1);
\draw[cube] (1,1,0) -- (1,1,1);
\draw[axis3] (0,0,0) -- (1.16,0,0) node[below right] {\(a_1\)};
\draw[axis3] (0,0,0) -- (0,1.15,0) node[left] {\(a_2\)};
\draw[axis3] (0,0,0) -- (0,0,1.15) node[above left=1pt] {\(a_3\)};
\coordinate (rA) at (0.38,0.58,0.42);
\coordinate (rB) at (0.50,0.40,0.28);
\coordinate (rS) at (0.88,0.98,0.70);
\draw[candA] (0,0,0) -- (rA) node[pos=1.10, above left=2pt, vlabel] {\(v_1\)};
\draw[candB] (0,0,0) -- (rB) node[pos=1.12, below right=2pt, vlabel] {\(v_2\)};
\draw[guide] (rA) -- (rS) -- (rB);
\draw[sumvec] (0,0,0) -- (rS) node[pos=0.98, right=4pt, vlabel] {\(v_1+v_2\)};
\draw[diffline] (rA) -- (rB) node[pos=0.55, left=8pt, above=3pt, difflabel] {\(v_1-v_2\)};
\end{scope}

\begin{scope}[xshift=5.05cm]
\node[font=\bfseries\scriptsize] at (0.48,0.50,1.48) {Complementary pair};
\fill[gray!5] (0,0,0) -- (1,0,0) -- (1,1,0) -- (0,1,0) -- cycle;
\fill[gray!9] (0,1,0) -- (1,1,0) -- (1,1,1) -- (0,1,1) -- cycle;
\fill[gray!13] (1,0,0) -- (1,1,0) -- (1,1,1) -- (1,0,1) -- cycle;
\draw[cube] (0,0,0) -- (1,0,0) -- (1,1,0) -- (0,1,0) -- cycle;
\draw[cube] (0,0,1) -- (1,0,1) -- (1,1,1) -- (0,1,1) -- cycle;
\draw[cube] (0,0,0) -- (0,0,1);
\draw[cube] (1,0,0) -- (1,0,1);
\draw[cube] (0,1,0) -- (0,1,1);
\draw[cube] (1,1,0) -- (1,1,1);
\draw[axis3] (0,0,0) -- (1.16,0,0) node[below right] {\(a_1\)};
\draw[axis3] (0,0,0) -- (0,1.15,0) node[left] {\(a_2\)};
\draw[axis3] (0,0,0) -- (0,0,1.15) node[above left=1pt] {\(a_3\)};
\coordinate (cA) at (0.72,0.20,0.24);
\coordinate (cB) at (0.16,0.72,0.65);
\coordinate (cS) at (0.88,0.92,0.89);
\draw[candA] (0,0,0) -- (cA) node[pos=1.08, below right=1pt, vlabel] {\(v_1\)};
\draw[candB] (0,0,0) -- (cB) node[pos=1.10, right=6pt, vlabel] {\(v_2\)};
\draw[guide] (cA) -- (cS) -- (cB);
\draw[sumvec] (0,0,0) -- (cS) node[pos=0.95, right=4pt, vlabel] {\(v_1+v_2\)};
\draw[diffline] (cA) -- (cB) node[pos=0.54, below=7pt, difflabel] {\(v_1-v_2\)};
\end{scope}
\end{tikzpicture}
}
\caption{Field-aware pair selection as a vector objective. Each axis is a query-field similarity component \(a_i\); actual retrieval vectors are \(m\)-dimensional. A redundant pair concentrates evidence in similar directions, so \(v_1-v_2\) is short. A complementary pair covers different field dimensions, increasing both the joint vector \(v_1+v_2\) and the difference vector \(v_1-v_2\) used by the pair score.}
\label{fig:field-aware-pair-selection}
\end{figure}

Demonstration diversity is not merely topical diversity. Two examples from different domains can still be redundant if both show easy scalar fields and neither illustrates nulls, arrays, or enum decisions. Conversely, two examples from a related family may be complementary if one illustrates absence and empty arrays while the other illustrates populated nested objects and label mapping. Tie-breaking favors stronger individual relevance and stable resource order, but the core decision is pair-level.

\subsection{Prompt Integration}

DRILLER integrates selected examples into the extraction prompt in one of two layouts. The general RAG layout is used when examples are related but not all same-schema matches. Each example is rendered as a self-contained demonstration with its instruction, schema context, source text, and expected output. This gives the model enough context to see how the example output follows from that example's schema before using it analogically.

The same-schema layout is used when all selected examples share the normalized target schema. The shared schema is rendered once, alongside the role instruction, extraction rules, root description when available, and guidance that examples are demonstrations rather than evidence. The examples then omit repeated schema blocks and focus on source-to-output behavior under the common structure. This reduces prompt length and makes contrasting field outcomes more salient.

Pipeline integration depends on the current output interface. In \texttt{Schema-RAG}, examples are rendered in direct final-value form, matching the object shape returned after generation. In \texttt{Reasoned-RAG}, examples are rendered with top-level \texttt{\{reasoning, value\}} objects because the temporary generation interface requires that structure. The reasoning strings demonstrate a local evidence protocol, while the \texttt{value} slots show the values that will later be unwrapped.

\texttt{Mixed-SC} uses the same retrieval machinery inside each trial. Direct trials use the \texttt{Schema-RAG} rendering, and reasoning-augmented trials use the \texttt{Reasoned-RAG} rendering. Each trial performs retrieval independently. Aggregation receives only postprocessed JSON candidates in the final schema shape.

Field-aware retrieval turns the local corpus into an inference-time schema interpretation aid. Same-schema retrieval exploits exact structural matches when available, field-level fallback finds type-compatible analogies, the pair objective favors complementary coverage, and prompt integration adapts examples to the direct, reasoned, and mixed pipelines.
